%Version 3 December 2023
% See section 11 of the User Manual for version history
%
%%%%%%%%%%%%%%%%%%%%%%%%%%%%%%%%%%%%%%%%%%%%%%%%%%%%%%%%%%%%%%%%%%%%%%
%%                                                                 %%
%% Please do not use \input{...} to include other tex files.       %%
%% Submit your LaTeX manuscript as one .tex document.              %%
%%                                                                 %%
%% All additional figures and files should be attached             %%
%% separately and not embedded in the \TeX\ document itself.       %%
%%                                                                 %%
%%%%%%%%%%%%%%%%%%%%%%%%%%%%%%%%%%%%%%%%%%%%%%%%%%%%%%%%%%%%%%%%%%%%%

%%\documentclass[referee,sn-basic]{sn-jnl}% referee option is meant for double line spacing

%%=======================================================%%
%% to print line numbers in the margin use lineno option %%
%%=======================================================%%

%%\documentclass[lineno,sn-basic]{sn-jnl}% Basic Springer Nature Reference Style/Chemistry Reference Style

%%======================================================%%
%% to compile with pdflatex/xelatex use pdflatex option %%
%%======================================================%%

%%\documentclass[pdflatex,sn-basic]{sn-jnl}% Basic Springer Nature Reference Style/Chemistry Reference Style


%%Note: the following reference styles support Namedate and Numbered referencing. By default the style follows the most common style. To switch between the options you can add or remove “Numbered” in the optional parenthesis. 
%%The option is available for: sn-basic.bst, sn-vancouver.bst, sn-chicago.bst%  
 
%%\documentclass[pdflatex,sn-nature]{sn-jnl}% Style for submissions to Nature Portfolio journals
%%\documentclass[pdflatex,sn-basic]{sn-jnl}% Basic Springer Nature Reference Style/Chemistry Reference Style
\documentclass[pdflatex,sn-mathphys-num]{sn-jnl}% Math and Physical Sciences Numbered Reference Style 
%%\documentclass[pdflatex,sn-mathphys-ay]{sn-jnl}% Math and Physical Sciences Author Year Reference Style
%%\documentclass[pdflatex,sn-aps]{sn-jnl}% American Physical Society (APS) Reference Style
%%\documentclass[pdflatex,sn-vancouver,Numbered]{sn-jnl}% Vancouver Reference Style
%%\documentclass[pdflatex,sn-apa]{sn-jnl}% APA Reference Style 
%%\documentclass[pdflatex,sn-chicago]{sn-jnl}% Chicago-based Humanities Reference Style

%%%% Standard Packages
%%<additional latex packages if required can be included here>

\usepackage{graphicx}%
\usepackage{multirow}%
\usepackage{amsmath,amssymb,amsfonts}%
\usepackage{amsthm}%
\usepackage{mathrsfs}%
\usepackage[title]{appendix}%
\usepackage{xcolor}%
\usepackage{textcomp}%
\usepackage{manyfoot}%
\usepackage{booktabs}%
\usepackage{algorithm}%
\usepackage{algorithmicx}%
\usepackage{algpseudocode}%
\usepackage{listings}%
\usepackage{hyperref}%
\usepackage{xr}

\externaldocument{supplementary}

%%%%

%%%%%=============================================================================%%%%
%%%%  Remarks: This template is provided to aid authors with the preparation
%%%%  of original research articles intended for submission to journals published 
%%%%  by Springer Nature. The guidance has been prepared in partnership with 
%%%%  production teams to conform to Springer Nature technical requirements. 
%%%%  Editorial and presentation requirements differ among journal portfolios and 
%%%%  research disciplines. You may find sections in this template are irrelevant 
%%%%  to your work and are empowered to omit any such section if allowed by the 
%%%%  journal you intend to submit to. The submission guidelines and policies 
%%%%  of the journal take precedence. A detailed User Manual is available in the 
%%%%  template package for technical guidance.
%%%%%=============================================================================%%%%

%% as per the requirement new theorem styles can be included as shown below
\theoremstyle{thmstyleone}%
\newtheorem{theorem}{Theorem}%  meant for continuous numbers
%%\newtheorem{theorem}{Theorem}[section]% meant for sectionwise numbers
%% optional argument [theorem] produces theorem numbering sequence instead of independent numbers for Proposition
\newtheorem{proposition}[theorem]{Proposition}% 
%%\newtheorem{proposition}{Proposition}% to get separate numbers for theorem and proposition etc.

\theoremstyle{thmstyletwo}%
\newtheorem{example}{Example}%
\newtheorem{remark}{Remark}%

\theoremstyle{thmstylethree}%
\newtheorem{definition}{Definition}%

\raggedbottom
%%\unnumbered% uncomment this for unnumbered level heads
\newcommand{\ekynprecision}{88.1}
\newcommand{\ekynrecall}{88.5}
\newcommand{\ekynfscore}{87.6}
\newcommand{\ekynfscorep}{78.2}
\newcommand{\ekynfscores}{92.3}
\newcommand{\ekynfscorew}{92.2}
\newcommand{\ekynprecisionp}{78.5}
\newcommand{\ekynprecisions}{94.1}
\newcommand{\ekynprecisionw}{91.8}
\newcommand{\ekynrecallp}{81.1}
\newcommand{\ekynrecalls}{91.2}
\newcommand{\ekynrecallw}{93.1}
\newcommand{\ekynprecisionstd}{4.3}
\newcommand{\ekynrecallstd}{6.4}
\newcommand{\ekynfscorestd}{5.6}
\newcommand{\ekynfscorepstd}{14.0}
\newcommand{\ekynfscoresstd}{5.0}
\newcommand{\ekynfscorewstd}{2.9}
\newcommand{\ekynprecisionpstd}{12.5}
\newcommand{\ekynprecisionsstd}{3.8}
\newcommand{\ekynprecisionwstd}{6.4}
\newcommand{\ekynrecallpstd}{16.6}
\newcommand{\ekynrecallsstd}{9.4}
\newcommand{\ekynrecallwstd}{3.8}

\newcommand{\liuprecision}{87.1}
\newcommand{\liurecall}{89.2}
\newcommand{\liuf}{88.1}
\newcommand{\liufp}{79.4}
\newcommand{\liufs}{92.5}
\newcommand{\liufw}{92.4}

\newcommand{\spindleprecision}{89.4}
\newcommand{\spindlerecall}{90.2}
\newcommand{\spindlefscore}{89.6}
\newcommand{\spindlefscorep}{83.6}
\newcommand{\spindlefscores}{92.1}
\newcommand{\spindlefscorew}{93.0}
\newcommand{\spindleprecisionp}{82.8}
\newcommand{\spindleprecisions}{89.9}
\newcommand{\spindleprecisionw}{95.5}
\newcommand{\spindlerecallp}{85.3}
\newcommand{\spindlerecalls}{94.5}
\newcommand{\spindlerecallw}{90.8}
\newcommand{\spindleprecisionstd}{2.6}
\newcommand{\spindlerecallstd}{3.9}
\newcommand{\spindlefscorestd}{2.9}
\newcommand{\spindlefscorepstd}{6.0}
\newcommand{\spindlefscoresstd}{2.8}
\newcommand{\spindlefscorewstd}{2.9}
\newcommand{\spindleprecisionpstd}{6.6}
\newcommand{\spindleprecisionsstd}{3.8}
\newcommand{\spindleprecisionwstd}{2.6}
\newcommand{\spindlerecallpstd}{9.6}
\newcommand{\spindlerecallsstd}{3.0}
\newcommand{\spindlerecallwstd}{4.3}

\begin{document}

\title[Article Title]{A deep learning software tool for automated sleep staging in rats via single channel EEG}

%%=============================================================%%
%% GivenName	-> \fnm{Joergen W.}
%% Particle	-> \spfx{van der} -> surname prefix
%% FamilyName	-> \sur{Ploeg}
%% Suffix	-> \sfx{IV}
%% \author*[1,2]{\fnm{Joergen W.} \spfx{van der} \sur{Ploeg} 
%%  \sfx{IV}}\email{iauthor@gmail.com}
%%=============================================================%%


\author[1]{\fnm{Andrew} \sur{Smith}}\email{andrewsmith@sc.edu}

\author[2]{\fnm{Snezana} \sur{Milosavljevic}}\email{snezana.milosavljevic@uscmed.sc.edu}

\author[2]{\fnm{Courtney} \sur{Wright}}\email{courtney.wright@uscmed.sc.edu}

\author[2]{\fnm{Charlie} \sur{Grant}}\email{charlie.grant@uscmed.sc.edu}

\author[2]{\fnm{Ana} \sur{Pocivavsek}}\email{ana.pocivavsek@uscmed.sc.edu}

\author*[1]{\fnm{Homayoun} \sur{Valafar}}\email{homayoun@cse.sc.edu}

\affil*[1]{\orgdiv{Computer Science and Engineering}, \orgname{University of South Carolina}, \orgaddress{\city{Columbia}, \postcode{29208}, \state{SC}, \country{USA}}}

\affil[2]{\orgdiv{Department of Pharmacology, Physiology, and Neuroscience}, \orgname{University of South Carolina School of Medicine}, \orgaddress{\city{Columbia}, \postcode{29208}, \state{SC}, \country{USA}}}

\abstract{Poor quality and poor duration of sleep have been associated with cognitive decline, diseases, and disorders. Therefore, sleep studies are imperative to recapitulate phenotypes associated with poor sleep quality and uncover mechanisms contributing to psychopathology. The composition of sleep in terms of sleep stages serves as a proxy for determining sleep quality. Currently, the gold standard for sleep staging is human inspection of polysomnography (PSG), which is time-consuming. To accelerate the analysis, we have developed a deep-learning-based software tool for automated sleep stage classification in rats. This study aimed to develop an automated method for classifying three sleep stages in rats (paradoxical sleep, slow-wave sleep, and wakefulness) using a deep learning approach based on single-channel EEG data. Single-channel EEG data were acquired from 16 rats, each undergoing two 24-hour recording sessions. The data were labeled by human experts in 10-second epochs corresponding to three stages: paradoxical sleep, slow-wave sleep, and wakefulness. A deep neural network (DNN) model was designed and trained to classify these stages using the raw temporal data from the EEG. The DNN achieved strong performance in predicting the three sleep stages, with an average F1 score of \ekynfscore\% over a cross-validated test set. The algorithm was able to predict key parameters of sleep architecture, including total bout duration, average bout duration, and number of bouts, with significant accuracy. Our deep learning model effectively automates the classification of sleep stages using single-channel EEG data in rats, reducing the need for labor-intensive manual annotation. This tool enables high-throughput sleep studies and may accelerate research into sleep-related pathologies. Furthermore, we provide over 700 hours of expert-scored sleep data, available for public use in future research studies.}

\keywords{sleep staging, deep learning, EEG, rats, sleep architecture, automated classification}

%%\pacs[JEL Classification]{D8, H51}

%%\pacs[MSC Classification]{35A01, 65L10, 65L12, 65L20, 65L70}

\maketitle
\section{Introduction}\label{introduction}
Sleep is essential for optimal health \cite{cirelli_is_2008}, yet more than one third of the global population reports problems with sleep \cite{jahrami_sleep_2021}. Preclinical sleep studies are imperative to recapitulate phenotypes associated with poor sleep (hyperarousal, cognitive impairment, slowed psychomotor vigilance, behavioral despair) and uncover mechanisms contributing to psychopathology \cite{graves_sleep_2003,rentschler_reducing_2024,siegel_sleep_2022,wright_stress_2023}.

Studies performed in rodents allow investigators to understand underlying mechanisms that drive sleep homeostasis \cite{rayan_sleep_2024}. Importantly, vigilance states that compose sleep are evolutionarily conserved amongst species, thereby rodent models can recapitulate the unique polysomnographic characteristics of sleep states \cite{rattenborg_evolution_2023}. The accurate determination of vigilance stages into wake, rapid eye movement (REM) sleep, and non-REM (NREM) is a critical step in analyzing acquired sleep-wake recordings. As REM and NREM sleep serve functions uniquely different from a wake state, correctly classifying vigilance states is imperative.

Currently, the gold standard for sleep staging (the determination of sleep stages) in rodents is by human inspection of polysomnography (PSG). PSG is a collection of signals obtained simultaneously during sleep (for example, electroencephalography (EEG), electromyography (EMG), electrooculography (EOG), etc.). Manual annotation of PSG by human experts, however, is a time-consuming process that severely limits efficiency. Further, human experts require extensive training that fails to eliminate interscorer and intra-scorer variability \cite{miladinovic_spindle_2019}. Therefore, developing a reliable, reproducible, and accurate automated method for rodent sleep staging is critical for more efficiently conducting preclinical sleep studies that are critical to improving our understanding of sleep and health outcomes.

Tethered rodents are combating confounds including limited mobility within recording cages and potential impacts on natural sleep states \cite{kramer_evaluation_2003,papazoglou_non-restraining_2016,rentschler_prenatal_2021}. For these reasons, it is advantageous to collect PSG via surgically implanted telemetry transmitters. Due to the size and nature of available telemetry transmitters for collecting PSG, the number of signals acquired simultaneously during sleep is limited by number of channels. Therefore, using a telemetry transmitter instead of tethering rodents via wires decreases the number of obtained signals available for sleep staging. Consistent with the principles of information theory, data collected from fewer signals significantly increases the complexity of downstream sleep staging \cite{shannon_mathematical_1948}.

Several published algorithms only require a few signals that may be collected via telemetry transmitters; however, they require a specific combination of EEG, EMG, and EOG recordings \cite{jha_slumbernet_2024,vallat_open-source_2021}. Therefore, developing an automated algorithm based on only a single-channel EEG signal is of particular interest because of its versatility across hardware and experimental settings. Further, automatic sleep staging based on a single-channel EEG signal collected via a telemetry transmitter would allow researchers the opportunity to investigate other signal types with remaining available channels, such as electrocardiogram (EKG).

Deep learning has revolutionized various fields, including time series analysis \cite{smith_toward_2024,chien_maeeg_2022,vaswani_attention_2023,he_deep_2015,krizhevsky_imagenet_2012}. Advances in deep learning have garnered interest in developing automated algorithms for sleep staging. Several such algorithms have been published in recent years. Several recent advancements in automated sleep staging are based on sleep studies in humans \cite{supratak_deepsleepnet_2017,vallat_open-source_2021,perslev_u-time_2019,Smith_Anand_Milosavljevic_Rentschler_Pocivavsek_Valafar_2022} wherein sleep staging standards differ significantly from those in rodent sleep studies \cite{rayan_sleep_2024}. Nearly all methods for automatic sleep staging rely on multi-channel EEG recordings and derived local field potential (LFP) signaling oscillations within the brain \cite{miladinovic_spindle_2019,gao_multiple_2016,gross_open-source_2009,stephenson_automated_2009,louis_design_2004,ellen_artificial_2021}. Of all the recently proposed deep learning algorithms for sleep staging, we have only come across two that have attempted automatic sleep staging in rodents via single-channel EEG \cite{tezuka_real-time_2021,liu_attention-based_2022}. 

In 2021, Tezuka et al. developed a CNN-LSTM neural network for sleep staging in mice from only a single-channel EEG; however, the authors use a private dataset, constrained to the light phase of the light-dark cycle, and include Zeitgeber Time (ZT) and the Discrete Fourier Transformation (DFT) as input to their model \cite{tezuka_real-time_2021}. The use of a private dataset limits researchers ability to confirm results or compare novel methods. Constraining data to the light phase of the light-dark cycle removes a dimension of variability that likely reduces the complexity of sleep staging \cite{shannon_mathematical_1948}. Including ZT and DFT as input may introduce bias and limit generalizability.

In 2022, Liu et al. developed a CNN in serial with an attention mechanism for sleep staging in rodents from only a single-channel EEG \cite{liu_attention-based_2022}. The authors use a public dataset from previously published work \cite{miladinovic_spindle_2019} comprised of 22 24-hour EEG recordings totaling 528 and provide the results of 22-fold leave-one-subject-out cross-validation \cite{liu_attention-based_2022}. Each of the 22 folds is a 24-hour EEG recording and, therefore, consists of both the light and the dark cycle. The authors, however, propose a complex preprocessing to their single-channel EEG signal before being input to their deep learning model including an array of hand-crafted time-frequency transformations.

Many existing sleep staging algorithms face challenges such as being proprietary, inadequately evaluated, dependent on complex data preprocessing, reliant on multiple EEG channels, or requiring additional signals such as EMG or EOG. Moreover, many are limited to human EEG data or trained solely on publicly available datasets. In contrast, we present an open-source, rigorously evaluated, end-to-end deep learning model trained on a novel dataset for sleep staging in rodents using single-channel EEG. Additionally, we contribute over 700 hours of expert-scored sleep data, providing a valuable resource for future research. We compare the performance of our model with closely related methods, including those by Tezuka et al. (2021) and Liu et al. (2022). Furthermore, using our newly introduced rodent dataset, termed SleepyRats, we demonstrate that our model accurately predicts key sleep architecture parameters, such as vigilance state durations, bout numbers, and average bout durations.

\section{Methods}\label{methods}
\subsection*{Data Acquisition Procedure}
Adult male and female Wistar rats (n=16 for training; n=14 for validation) were used in experiments in a facility fully accredited by the Association for Assessment and Accreditation of Laboratory Animal Care (AAALAC). Rats were kept on a 12 h/12 h light-dark cycle. All protocols were approved by the Institutional Animal Care and Use Committee (IACUC) at the University of South Carolina and were in accordance with the National Institutes of Health Guide for the Care and Use of Laboratory Animals.

Rats were implanted with EEG/EMG telemetry transmitters (PhysioTel HD-S02, Data Sciences International (DSI), St. Paul, MN, USA). Briefly, under isoflurane anesthesia, rats were placed in a stereotaxic frame (Stoelting Co., Wood Dale, IL, USA). The telemetry transmitter was implanted intraperitoneally through a dorsal incision of the abdominal region. After an incision at the midline of the head was made, two EEG leads were secured to two surgical stainless-steel screws (P1 Technologies, Roanoke, VA, USA) inserted into 0.5-mm burr holes at 2.0 mm anterior/+1.5 mm lateral and 7.0 mm posterior/-1.5 mm lateral relative to bregma. The two EEG leads were secured using acrylic dental cement (Stoelting Co., Wood Dale, IL, USA). Two EMG leads were inserted into the dorsal cervical neck muscle about 1.0 mm apart and sutured into place. The skin along the head was sutured, and animals were recovered for a minimum of 7 days before sleep data acquisition.

Data used in training our model were obtained from rats that were subject to 2 discrete sleep recording sessions each consisting of exactly 24 hours \cite{milosavljevic_kynurenine_2023}. Sleep data were acquired in a quiet, designated room where rats remained undisturbed for the sleep recording duration on a PC running Ponemah 6.10 software (DSI). EEG/EMG data were acquired with a sampling rate of 500 Hz. Digitized signal data were exported for both recording sessions in European Data Format (EDF), a standard file type for raw PSG. An expert in sleep staging manually annotated each recording with sleep stages in 10-second epochs offline using NeuroScore 3.0 software (DSI).

For each of the 2 recording sessions, 16 discrete 24-hour EEG/EMG signals resulted from data collection (one for each rat). We group the pairs of 24-hour PSG recordings for each rat into a 48-hour PSG recording and ensure that data is processed in a temporally consistent manner. It is however critical that we maintain rat identity throughout the entire investigation for proper training and evaluation of our algorithm. Therefore, we refer to the shape of the dataset as 16 rats x 48-hour PSG recording, as shown in Fig.~\ref{dataset}.

\subsection*{Validation}
The validation paradigm is critical for proper model evaluation and establishing robust results \cite{Goodfellow-et-al-2016}. The goal of defining a validation paradigm is to simulate the use of a model deployed in the real world. Most machine learning applications aim to develop a system capable of generalization as opposed to memorization. Generalization refers to a model's ability to perform well when provided new data that it was not explicitly trained on, whereas memorization refers to a model's tendency to excessively rely on specific details and examples from the training data, failing to capture underlying patterns or relationships. Because of a deep learning model's propensity to overfit to training data, leakage of information between training and testing sets may result in overoptimistic results. Subject-based datasets may be prone to leakage of information due to systematic variation between subjects (exposure of a network to any given subject in both training and testing). To remove the possibility of information leakage, we use a cross-validation scheme (a particular form of validation in general) where a given rat's PSG data is contained within only one-fold. This is referred to as leave-one-out cross-validation, a special case of cross-validation, where the number of folds is equal to the number of groups. Since there are 16 rats in our dataset, we define each rat (consisting of a 48-hour PSG recording) as a group. Therefore, for each fold, we train a model on the 15 others and evaluate the last rat. This approach typically provides the most reliable validation of generalization under the condition of data scarcity \cite{Hastie_Tibshirani_Friedman_2009}.

\subsection*{Model}
We designed our model to independently encode EEG epochs into information-rich low-dimensional vectors then incorporate sequential information from neighboring epochs to make a single sleep stage prediction. To accomodate this design, we combine a convolutional neural network \cite{krizhevsky_imagenet_2012} with a Long-Short Term Memory neural network \cite{hochreiter_long_1997}. We designed our model to take a sequence of 9 10-second epochs as input and precit the sleep stage of the center (5th) epoch based on the common practice of our expert labelers; however, one could easily adapt the architecture to take any sequence length. Our model architecture is shown in Fig.~\ref{model}.

\begin{figure}[H]
    \centering
    \includegraphics[width=\textwidth]{fig1.pdf}
    \caption{
        \textbf{Model diagram of our algorithm.} The model accepts a sequence of 9 10-second epochs to predict the class of the center 10-second epoch (epoch i, as denoted in the figure) Each 10-second epoch in the sequence is independently encoded and then forwarded to a Long Short-Term Memory (LSTM) neural network to share information between neighboring 10-second epochs. The 64-dimensional LSTM output vector is linearly projected down to three dimensions and mapped to a probability distribution over 3 values (corresponding to 3 vigilant states) by the Softmax function.
    }\label{model}
\end{figure}

\noindent \textbf{Epoch Encoder:}\quad During the first stage of the network, each 10-second epoch in the input sequence is independently encoded by the epoch encoder, shown in Fig.~\ref{model}. The epoch encoder is a RegNet \cite{radosavovic_designing_2020} modern variant of the convolutional neural network \cite{krizhevsky_imagenet_2012} that investigated optimal design of Residual Neural Networks \cite{he_deep_2015}. The epoch encoder consists of a stack of 3 identical residual blocks \cite{he_deep_2015} and a final global average pooling \cite{lin_network_2014}, except that the number of convolution channels vary. The number of convolution channels is shown within each residual block in Fig.~\ref{model_components}. Each residual block consists of 3 sequences of convolution, layer normalization \cite{ba_layer_2016}, and ReLU nonlinearity \cite{fukushima_visual_1969}, as well as an elementwise residual connection before the final ReLU. The convolution kernels for each residual block are size 8, 5, and 3, in that order.

\noindent \textbf{LSTM and Classifier Head:}\quad The resultant sequence of spatially encoded feature maps is forwarded to a temporally aware neural network, resulting in a single output feature map that integrates spatial and temporal information. Lastly, this spatio-temporally-encoded feature map is presented to a final classifier that maps 3 output states onto a probability distribution, each of which corresponds to one of 3 vigilant states. Additional details regarding the network architecture can be found in the supplementary material Fig.~\ref{model_components}.

\subsection*{Preparing Dataset For Training}
We train the model to accept an input sequence of 9 10-second epochs and output a probability distribution over 3 classes corresponding to 3 stages (paradoxical sleep, slow-wave sleep, and wakefulness) for the 10-second epoch in the center of the sequence. To allow the evaluation of the first 4 10-second epochs, we pad each 1-dimensional EEG signal (one 24-hour PSG recording) with 4 10-second epochs with zero value on each end. Each zero-padded PSG signal is partitioned into 10-second epochs and transformed using a simple moving window of 9 10-second epochs with a stride of 1. Initial expert sleep staging remains unchanged; therefore, each window of 9 sequential 10-second epochs is annotated by the sleep stage of the epoch in the center of the window as desired. We take a 5\% of the training set on each fold for validation.

\subsection*{Optimization}
We trained our model using mini batches of size 32 epochs (input sequences of 9 10-second epochs and the corresponding sleep staging) until validation loss did not improve for 30 epochs. We perform hyperparameter optimization using the validation loss. During training, we used cross-entropy \cite{kullback_information_1951} to compare model predictions to reference sleep staging. For model weight optimization, we used Adam \cite{kingma_adam_2017}. All the pre and post data processing steps were completed in Python using PyTorch as the sole machine learning platform.

\subsection*{Evaluation Metrics}
Defining good evaluation metrics is crticial for establishing proper results in deep learning \cite{Goodfellow-et-al-2016}. Sleep staging is inherently an imbalanced classification problem. Using accuracy as a performance metric in an imbalanced classification problem provides overoptimistic results. A good performance metric for evaluating an imbalanced classification is F1-score, the harmonic mean between precision and recall. In a classification setting with multiple classes, we define the F1 Score for a given class $C$ as
\begin{gather*}
\text{F1 Score}_C = \frac{2 (\text{precision}_C)(\text{recall}_C)}{\text{precision}_C + \text{recall}_C},
\text{precision} = \frac{\text{TP}_C}{\text{TP}_C+\text{FP}_C},\text{recall} = \frac{\text{TP}_C}{\text{TP}_C+\text{FN}_C}
\end{gather*}
where TP is the number of true positives, FP the number of false positives, and FN the number of false negatives. We define the macro F1 Score as the average F1 Score over each class
$$
\frac{\Sigma_C \text{F1 Score}_C}{N}
$$
where there are $N$ classes.
\subsection*{Comparison Against Previous Work}
In addition to testing the developed system using the procedure stated in section 3.2, we further tested the performance of our neural network against a recently developed sleep-staging algorithm published by Liu et al.\cite{liu_attention-based_2022}. Their algorithm is based on a deep learning approach that was developed and evaluated using a publicly available dataset \cite{miladinovic_spindle_2019}. We evaluated our algorithm using the same dataset to provide a direct comparison of performance and to establish the generalization capabilities of our trained network.

The validation scheme for our dataset is comparable to the validation approach in Liu et al.\cite{liu_attention-based_2022}, wherein a leave-one-out cross-validation over 22 folds, each corresponding to a 24-hour sleep recording from a unique rodent, was utilized. The SPINDLE dataset is composed of multiple cohorts of animals sampled at varying sampling rates (128 Hz, 200 Hz, and 512 Hz) \cite{miladinovic_spindle_2019}. To adapt the SPINDLE dataset to our model, we resampled each EEG recording to 500 Hz using the Fourier Transform method and performed the same zero-padding technique described above (see 3.4 Data Preparation). Additionally, the SPINDLE dataset defines the length of an epoch as 4 seconds and performs the classification at that resolution. Our model outputs sleep staging at the 10-second resolution; therefore, our predictions are not directly comparable to the SPINDLE reference. To alleviate this problem, we resampled our prediction signal from the 10-second resolution to the 4-second resolution (i.e. up sampling). By up sampling our prediction signal (as opposed to down sampling the SPINDLE reference signal), no information is lost from the SPINDLE reference signal and no information is gained in the prediction signal; therefore, the problem's complexity remains unchanged. Up sampling a signal from 0.1 Hz to 0.25 Hz is non-trivial since 0.1 does not divide 0.25, resulting in 1 unaligned reference 4-second epoch for every 2 predictive 10-second epochs as shown in Fig. 4. For every 2 consecutive 10-second epochs in our prediction signal, there are 5 consecutive 4-second epochs in the SPINDLE reference signal. The first two and last two reference 4-second epochs align within 10-second epochs in our prediction signal and are directly comparable. However, every third 4-second epoch in the SPINDLE reference signal spans 2 10-second epochs in our prediction signal. To account for this, we compare this SPINDLE reference 4-second epoch to both predictive 10-second epochs which it spans and take the mean of the metric being computed. For example, the third 4-second epoch in the SPINDLE reference signal (P) spans 2 10-second epochs in our prediction signal (P and S, respectively); therefore, this epoch contributes  5 to accuracy (Fig. \ref{spindle}A).

\subsection*{Validation of conclusions on vigilance state duration and architecture}
Data used in validating our model for conclusions on vigilance state duration and architecture were obtained from rats that were subject to discrete sleep recording sessions, as described above. An expert in sleep staging manually annotated each recording with sleep stages in 10-second epochs offline using NeuroScore 3.0 software (DSI) for the first 4-hour PSG recording. Comparisons were made between LSTM scoring the human scoring in 1-hour bins from ZT 0 to 4 evaluating vigilance state duration, bout number, and average duration of each bout. Comparisons were made between LSTM scoring and human scoring using a two-way repeated measures (RM) analysis of variance (ANOVA) with ZT and scoring as within-subject factors. Post hoc analysis was conducted to evaluate changes across time with Dunnett test compared to ZT 0. All statistical analyses were performed using Prism 9.0 (GraphPad Software, La Jolla, CA, USA) and significance was defined as P < 0.05.

\subsection*{Data and Code Availability}
All data can be requested from our website (\href{https://ifestos.cse.sc.edu/datasets/ekyn.tar.gz}{https://ifestos.cse.sc.edu/datasets/ekyn.tar.gz}). The SPINDLE dataset can be found at \href{https://sleeplearning.ethz.ch/paper/}{https://sleeplearning.ethz.ch/paper/}. The source code and documentation of the algorithm are available at \href{https://github.com/smithandrewk/ml-sleep}{https://github.com/smithandrewk/ml-sleep}. The Python code to reproduce all the results and figures of this paper can be found at \href{https://github.com/smithandrewk/ml-sleep}{https://github.com/smithandrewk/ml-sleep}. All analyses were conducted in Python 3.11 using PyTorch 1.13.1.

\section{Results}\label{results}
\subsection*{KYNA Dataset}
The dataset consisted of 32 unique 24-hour PSG recordings obtained from a cohort of 16 rats subjected to 2 discrete sleep recording sessions. We trained our model  with fixed architecture and hyperparameters over all 16 folds (where each fold corresponds to an individual rat). The training set, for each of the 16 folds, consisted of 30 24-hour EEG signals (259200 10-second epochs). A small portion of the training set (5\%) was taken to implement early stopping for model training . Early stopping is a regularization technique that monitors the model's performance on a validation set and halts training once performance on the validation set starts to degrade, thereby preventing overfitting and improving generalization to unseen data. The testing set, for each of the 16 folds, consisted of 2 24-hour EEG signals (17280 10-second epochs). The overall performance of the algorithm is described in Fig.~\ref{ekyn_grid}A. The box plot shows the distribution of performance over 16 folds of cross-validation. The mean f1-score, calculated across all 16 folds from leave-one-out cross-validation, was \ekynfscore\%. One outlier incurred several misclassifications, resulting in a recall of approximately 78\% and an f1-score of approximately 83\% (Fig.~\ref{ekyn_grid}A).

\begin{figure}[H]
    \centering
    \includegraphics[width=\textwidth]{fig2.pdf}
    \caption{
        \textbf{Performance of the algorithm on our dataset created for the model training over 16 folds of cross-validation.}\quad (A) Precision, recall, and f1-score over 16 folds of cross-validation are shown by a boxplot. (B) Dimensionality reduction of EEG encodings colored by sleep stage. Our network encodes EEG signals to be well-separated and therefore serves as the basis for accurate sleep staging. (C) Confusion matrix. The diagonal elements represent the percentage of 10-second epochs that were correctly classified by the algorithm (recall), whereas the off-diagonal elements show the percentage of 10-second epochs mislabeled by the algorithm. (D) Duration of each sleep stage for predicted signal and reference signal over 16 folds of cross-validation expressed as a proportion of the total duration of the given EEG recording.  Pairs of boxplots are shown for each stage where the left box depicts the predicted distribution, and the right box depicts the reference distribution. P - paradoxical sleep, S - slow-wave sleep, W - wakefulness.
    }\label{ekyn_grid}
\end{figure}

\begin{table}[ht]
    \caption{Performance of our model on our dataset broken down by class. Metrics are the mean and standard deviation over 16 folds of cross-validation for f1-score, recall, and precision.}
    \label{kyna_table}
    \vskip 0.15in
    \begin{centering}
    \begin{small}
    \begin{sc}
    \begin{tabular}{lcccr}
    \toprule
    Class & F1 & recall & precision \\
    \midrule
    P & \ekynfscorep $\pm$ \ekynfscorepstd & \ekynrecallp $\pm$ \ekynrecallpstd & \ekynprecisionp $\pm$ \ekynprecisionpstd\\
    S & \ekynfscores $\pm$ \ekynfscoresstd & \ekynrecalls $\pm$ \ekynrecallsstd & \ekynprecisions $\pm$ \ekynprecisionsstd\\
    W & \ekynfscorew $\pm$ \ekynfscorewstd & \ekynrecallw $\pm$ \ekynrecallwstd & \ekynprecisionw $\pm$ \ekynprecisionwstd\\
    \bottomrule
    \end{tabular}
    \end{sc}
    \end{small}
    \end{centering}
    \vskip -0.1in
\end{table}

The model compresses 45,000-dimensional  (500 Hz * 10 Seconds * 9 epochs) raw EEG inputs into a 64-dimensional  latent spatiotemporal representation before being linearly projected into three dimensions  for sleep staging. To visualize similarities between these 64-dimensional feature maps, we show the results of tSNE \cite{maaten_visualizing_2008}, a dimension reduction algorithm designed to preserve high-dimensional neighborhoods in Fig.~\ref{ekyn_grid}B. As expected, slow-wave sleep and wakefulness stages are well-clustered whereas many wake and paradoxical sleep stages' embeddings escape into the other's neighborhood. The similarity between EEG waveforms during wakefulness and paradoxical sleep almost certainly contributes to these mixtures.

We tested the sleep staging performance of individual vigilant states (paradoxical sleep, slow-wave sleep, and wakefulness) by computing a recall confusion matrix (Fig.~\ref{ekyn_grid}C). The average recall for paradoxical sleep is \ekynrecallp\%, and most misclassifications were classified as wakefulness. Both slow-wave sleep and wakefulness have a recall over \ekynrecalls\%. Class specific recall, precision, and f1 score can be found in Table~\ref{kyna_table}.

The dynamics of the composition of sleep, in terms of sleep stages, are complex but important for interpretations of sleep quality in preclinical studies. A few important parameters of sleep composition, beyond merely total duration, can thus serve as a unique characteristic for evaluation of model performance and include total bout duration , average bout duration, and number of bouts (all distributed by vigilant state). We investigated these parameters (Fig.~\ref{proxy_grid}) and show the total duration of each stage as a proportion of EEG recording time in Fig.~\ref{ekyn_grid}D.  

Tezuka et. al \cite{tezuka_real-time_2021} report a mean f1-score of 80.8 across 10 folds of cross-validation, which is notably lower than the F1-score achieved by our model. Consequently, we focus our comparison on SPINDLE, where the performance gap is more meaningful for further evaluation.

\subsection*{SPINDLE Dataset}
It is critical to emphasize that our model was optimized for sleep staging with 10-second epochs and that stages in the SPINDLE dataset were scored with 4-second epochs. In evaluating our model on the SPINDLE dataset, we find that, even by unique characteristic of 10-second epochs, we outperform the algorithm proposed by Liu et al. \cite{liu_attention-based_2022} despite theirs being optimized specifically for the SPINDLE dataset .
The results of 22 folds leave-one-out cross-validation on the SPINDLE dataset are reported in Table~\ref{metrics_spindle}. Notably, using the same dataset and validation scheme, our model achieves a mean f1-score of \spindlefscore\% compared to \liuf\% in Liu et al. \cite{liu_attention-based_2022}. Though we achieve higher performance in all three classes, most of the performance increase comes from the paradoxical sleep class where our model achieves f1-score of \spindlefscorep\% as opposed to \liufp\% in Liu et al. \cite{liu_attention-based_2022}. It is critical to emphasize that our entry in Table~\ref{metrics_spindle} reflects our performance after up sampling our model's 10-second prediction signal to be a 4-second prediction signal.
We performed further analysis on the distribution of performance across the 22 folds of cross-validation, the results of which are shown in Figure \ref{spindle_grid}.

Notably, one outlier achieves an f1-score of approximately 83\% . To further understand model performance broken down by class, we present the mean recall confusion matrix over the 22 folds of cross-validation  (Fig.~\ref{spindle_grid}C). Most misclassifications oc slackcur between paradoxical sleep and wakefulness, as expected due to the similarity between their EEG waveforms.
For each input sample, our model outputs a probability distribution over all possible stages; therefore, we can quantify the model's confidence in predicting a certain class. For each fold, we take the average confidence of the model and compare this to the macro f1-score  for that recording (Fig.~\ref{spindle_grid}B). Linear regression between these two variables suggests a positive relationship between a model's average confidence and the model's performance.
As presented for the dataset created for the model training, we investigated the proportion of recording time that each stage constituted as well as the proportion of total EEG recording time predicted by the model (Fig.~\ref{spindle_grid}D). Though for each class and each parameter the distribution of the SPINDLE reference values and predicted values is comparable, it is still unclear how well the model performs for each EEG recording. Thus, we perform linear regression between the SPINDLE reference values and predicted values for all combinations of classes and all three parameters (Fig.~\ref{proxy_grid}). Strong positive linear correlation was confirmed for all combinations of classes and parameters. If the model was a perfect predictor, all individual data points would fall on the x = y line.
A common method of visualizing sleep stages for a given EEG recording is a hypnogram. The reference hypnogram is shown in Fig.~\ref{hypnogram}A and our model's predicted hypnogram in shown in Fig.~\ref{hypnogram}B. Further, the EEG signal that constitutes the only input to the model, as it is an end-to-end deep learning model, is shown in Fig.~\ref{hypnogram}C.
Lastly, the predicted probability distribution over all possible stages by the model for the EEG signal is shown in Fig.~\ref{hypnogram}D. We take the highest probability as the predicted class, and the minimum confidence  with which our model could predict a certain class is slightly above 0.33. Hypothetically, higher average confidence values result in higher network performance. We notice that the confidence level of the network fluctuates throughout the recording; in longer bouts, confidence seems to increase, and in shorter bouts, confidence seems to decrease. Region from epoch 150 to epoch 250 on Fig.~\ref{hypnogram}D shows where the network is almost 100\% confident.

\subsection*{SleepyRat Dataset}
SleepyRat study was conducted in a cohort of rats to evaluate our novel sleep scoring network. We evaluated vigilance state durations and architecture comparing the findings of LSTM scored date to an expert human. Taken together, we determined no significant differences in the data due to scoring (Table~\ref{architecture_table}), but found significant impact of time of day, ZT, on vigilance state duration and architecture (bout number and average duration of vigilance bouts) parameters. In the parameters that a main effect of ZT was determined, post hoc Dunnett test analysis revealed significant differences between ZT 0-1 and subsequent time bins. Of note, findings between LTSM scoring and human scoring on the impact of time on vigilance state duration and architecture parameters aligned (Figure~\ref*{architecture}).

\begin{table}[ht]\label{architecture_table}
    \centering
    \caption{\textbf{Results of SleepyRat dataset statistical analysis comparing LSTM scoring and human expert scoring.}\quad P values are indicated for main effects of scoring, ZT, and interaction between scoring and ZT. * indicates significant effect wherein P $<$ 0.05. P values of two-way RM ANOVA.}
    \begin{tabular}{lccc}
    \toprule
    \textbf{Vigilance State} & \textbf{Scoring} & \textbf{ZT} & \textbf{Interaction} \\
    \midrule
    \multicolumn{4}{l}{\textbf{Duration}} \\
    S, NREM & 0.4852 & 0.0423* & 0.3153 \\
    P, REM & 0.1281 & $<$0.0001* & 0.2553 \\
    W, Wake & 0.3713 & 0.0008* & 0.2074 \\
    \midrule
    \multicolumn{4}{l}{\textbf{Bout Number}} \\
    S, NREM & 0.9327 & 0.0341* & 0.0516 \\
    P, REM & 0.1857 & $<$0.0001* & 0.4208 \\
    W, Wake & 0.8179 & 0.0535 & 0.0876 \\
    \midrule
    \multicolumn{4}{l}{\textbf{Average Bout Duration}} \\
    S, NREM & 0.4614 & 0.1865 & 0.2871 \\
    P, REM & 0.3060 & 0.4495 & 0.3416 \\
    W, Wake & 0.4412 & 0.0023* & 0.0624 \\
    \bottomrule
    \end{tabular}
    \end{table}

\begin{figure}[H]
    \centering
    \includegraphics[width=\textwidth]{fig3.pdf}
    \caption{
        \textbf{Sleep-wake duration and architecture data do not significantly differ between LSTM and human expert scoring of SleepyRat dataset. }\quad (A) NREM duration, (B) REM duration, (C) Wake duration, (D) NREM bout number, (E) REM bout number, (F) Wake bout number, (G) average NREM bout duration, (H) average REM bout duration, (I) average Wake bout duration. Two-way RM ANOVA with post-hoc Dunnett test. *  P < 0.05, ** P < 0.01, *** P < 0.001, **** P < 0.0001. Data are mean $\pm$ SEM. N =14 per group.
    }\label{architecture}
\end{figure}

\section{Discussion}\label{discussion}
The current gold standard for sleep staging involves a trained human expert manually annotating polysomnography (PSG), typically using EEG and EMG signals, to classify sleep stages. Developing a reliable, automated algorithm for sleep staging is critical for improving the efficiency of EEG signal processing. An automated method that uses only single-channel EEG recordings would provide a significant advancement by reducing the need for multi-channel recordings and manual annotations.

Human experts, after extensive training, rely on visual inspection of signals (such as EEG and EMG) and feature-based tools (like power-spectral density) to classify sleep epochs into vigilant stages. The patterns within these signals (e.g., amplitude and frequency) ultimately guide expert classifications. To replicate this process, we implemented a convolutional neural network to encode visual patterns contained in raw EEG signals down to a representation more amenable to classification. Each 10-second epoch is processed independently, resulting in a set of feature maps.

To capture temporal relationships between EEG patterns, we introduced a "temporal encoder" following the spatial encoder. This temporal encoder is designed to share information across time by considering 4 preceding and 4 succeeding epochs (for a total window of 9 epochs, or 90 seconds) centered around the target epoch to be classified. This mirrors the 90-second window used by human scorers when analyzing PSG. The temporal encoder is based on a Long Short-Term Memory (LSTM) network, which is capable of learning and integrating temporal dependencies from the data. The combination of spatial and temporal encoders allows our model to produce an intermediate representation rich in spatiotemporal information.

Not only does our model effectively capture spatiotemporal information, but it also ensures that the representations of different sleep stages are well-separated in the latent space, which is key for accurate classification. This capability of deep learning—to automatically transform raw data into meaningful representations for downstream tasks—highlights the strength of our approach. Our end-to-end model takes raw single-channel EEG signals as input and outputs sleep stage classifications, avoiding the need for feature extraction steps that could introduce bias or redundancy.

To further validate the usability of our model, we conducted an extensive evaluation using the SleepyRats dataset. This dataset consists of over 700 hours of single-channel EEG recordings from a cohort of 16 rats recorded across two 24-hour sleep sessions. The dataset provides a robust foundation for evaluating the model's performance in terms of vigilance state durations, bout numbers, and the average duration of bouts. A two-way repeated measures ANOVA was performed to compare the LSTM-based sleep staging results with human expert scoring. No significant differences were found between the LSTM predictions and human annotations for sleep stage durations and architecture parameters, confirming the model's reliability for preclinical sleep analysis. As shown in Figure \ref{architecture}, the consistency in the number of bouts and the average bout durations across different sleep stages between LSTM scoring and human scoring further validates the model's effectiveness.

Additionally, we provide comparisons to existing datasets like SPINDLE, which consisted of 22 folds and under 500 hours of EEG recordings. Proper evaluation of a deep learning model becomes the most critical step in the development of machine learning systems. Our model was rigorously evaluated using leave-one-out cross-validation on both the newly introduced SleepyRats dataset and the SPINDLE dataset. In total, 38 deep learning models were developed during the validation phase, and performance was measured using F1 score, precision, and recall.

Through the introduction of the SleepyRats dataset and the validation performed using ANOVA, we have demonstrated that our deep learning model not only automates sleep staging effectively but also provides reliable estimates of sleep architecture parameters. This validation, combined with the model's performance on cross-validation, ensures its usability for large-scale preclinical sleep analysis, offering a significant reduction in the time and resources required for manual sleep staging.
\section*{Acknowledgements}
Thanks to reviewers who gave useful comments, to colleagues who contributed to the ideas, and to funding agencies that provided financial support.
\section*{Author Contributions}
Andrew Smith contributed to the conception and design of the study, data analysis and interpretation, and drafting the manuscript. Courtney Wright contributed to data acquisition and provided revisions. Charlie Grant contributed to data analysis and revisions. Snezana Milosavljevic contributed to data acquisition and revisions. Ana Pocivavsek contributed to the conception and design of the study, provided revisions, and served as an advisor. Homayoun Valafar contributed to the conception and design of the study, provided revisions, and served as an advisor. All authors reviewed and approved the final manuscript, and agree to be accountable for all aspects of the work.
\section*{Funding}
Research reported in this publication was supported by funds awarded to Dr. Valafar from the National Institute of General Medical Sciences of the National Institutes of Health under Award Number P20GM103499. Research reported in this public was also supported by R01 NS 102209 and R21 AG080335 awarded to Dr. Ana Pocivavsek.
\section*{Competing Interests}
The authors have nothing to disclose.

\bibliography{sn-bibliography}% common bib file
%% if required, the content of .bbl file can be included here once bbl is generated
%%\input sn-article.bbl
\section*{Figure Legends}
\noindent
\textbf{Figure 1. Model diagram of our algorithm.}\quad The model accepts a sequence of 9 10-second epochs to predict the class of the center 10-second epoch (epoch i, as denoted in the figure) Each 10-second epoch in the sequence is independently encoded and then forwarded to a Long Short-Term Memory (LSTM) neural network to share information between neighboring 10-second epochs. The 64-dimensional LSTM output vector is linearly projected down to three dimensions and mapped to a probability distribution over 3 values (corresponding to 3 vigilant states) by the Softmax function.
\vspace{12pt} % Adds space between legends

\noindent
\textbf{Figure 2. Performance of the algorithm on our dataset created for the model training over 16 folds of cross-validation.}\quad (A) Precision, recall, and f1-score over 16 folds of cross-validation are shown by a boxplot. (B) Dimensionality reduction of EEG encodings colored by sleep stage. Our network encodes EEG signals to be well-separated and therefore serves as the basis for accurate sleep staging. (C) Confusion matrix. The diagonal elements represent the percentage of 10-second epochs that were correctly classified by the algorithm (recall), whereas the off-diagonal elements show the percentage of 10-second epochs mislabeled by the algorithm. (D) Duration of each sleep stage for predicted signal and reference signal over 16 folds of cross-validation expressed as a proportion of the total duration of the given EEG recording.  Pairs of boxplots are shown for each stage where the left box depicts the predicted distribution, and the right box depicts the reference distribution. P - paradoxical sleep, S - slow-wave sleep, W - wakefulness.
\vspace{12pt} % Adds space between legends

\noindent
\textbf{Figure 3. Sleep-wake duration and architecture data do not significantly differ between LSTM and human expert scoring of SleepyRat dataset.}\quad (A) NREM duration, (B) REM duration, (C) Wake duration, (D) NREM bout number, (E) REM bout number, (F) Wake bout number, (G) average NREM bout duration, (H) average REM bout duration, (I) average Wake bout duration. Two-way RM ANOVA with post-hoc Dunnett test. *  P < 0.05, ** P < 0.01, *** P < 0.001, **** P < 0.0001. Data are mean $\pm$ SEM. N =14 per group.
\end{document}
